% Regression test for \catName's auto-detection heuristic, tightened
% 2026-09-01 (see semantic-markup.dtx, the "Tightened" correction
% under "cat vs. catname (resolved)"). A single-token test alone
% mis-classified real arguments from the source papers -- decorated
% symbols (\catname{C^2}, \catname{E^i}) and a two-letter product-
% category symbol (\catname{XY}, 44 real call sites) -- as bold
% upright rather than script, because none of them is a single token
% even though the papers call all three via \catname (the "arbitrary
% symbol" choice), never \cat (the "long name" choice).
%
% Verifies, for every distinct argument actually found in
% LCS.arXiv.V2.tex/M-atlas.tex's \catname/\cat call sites (17 unique
% arguments across all 3413 calls), that \catName reproduces the
% papers' own existing choice: script for everything called via
% \catname, bold upright for everything called via \cat. Fails hard
% (raises a LaTeX error) on any mismatch rather than requiring visual
% inspection, since this is exactly the kind of silent typeface flip
% a future edit to the heuristic could reintroduce.
\documentclass{article}
\usepackage{semantic-markup}

\ExplSyntaxOn

\let\reservedmathscr\mathscr
\let\reservedmathbf\mathbf
\tl_new:N \l_lsmtest_got_tl
\renewcommand{\mathscr}[1]{\tl_gset:Nn \l_lsmtest_got_tl {script} \reservedmathscr{#1}}
\renewcommand{\mathbf}[1]{\tl_gset:Nn \l_lsmtest_got_tl {bold} \reservedmathbf{#1}}

% #1 = argument, #2 = expected (script|bold)
\cs_new_protected:Npn \lsmtestCatName #1#2
  {
    \tl_clear:N \l_lsmtest_got_tl
    \hbox{$\catName{#1}$}
    \tl_if_eq:NnTF \l_lsmtest_got_tl {#2}
      { \iow_term:x { ok~~catName\{ \tl_to_str:n{#1} \}~~->~~#2 } }
      {
        \iow_term:x
          { FAIL~~catName\{ \tl_to_str:n{#1} \}~~expected~#2~~got~
            \l_lsmtest_got_tl }
        \msg_error:nnxx { lsmtest } { mismatch }
          { \tl_to_str:n{#1} } {#2}
      }
  }
\msg_new:nnn { lsmtest } { mismatch }
  { catName\{#1\}~expected~#2~but~did~not~get~it. }

\ExplSyntaxOff

\begin{document}

% The 12 plain single-letter arguments actually used (called via
% \catname in the source papers -- expect script).
\lsmtestCatName{A}{script}
\lsmtestCatName{B}{script}
\lsmtestCatName{C}{script}
\lsmtestCatName{E}{script}
\lsmtestCatName{G}{script}
\lsmtestCatName{M}{script}
\lsmtestCatName{S}{script}
\lsmtestCatName{T}{script}
\lsmtestCatName{U}{script}
\lsmtestCatName{V}{script}
\lsmtestCatName{X}{script}
\lsmtestCatName{Y}{script}

% The two long, conventional category names (called via \cat --
% expect bold upright).
\lsmtestCatName{Set}{bold}
\lsmtestCatName{Top}{bold}

% The three decorated/concatenated symbols that a single-token-only
% test gets wrong (called via \catname despite not being single
% tokens -- expect script; this is the actual regression being
% guarded).
\lsmtestCatName{C^2}{script}
\lsmtestCatName{E^i}{script}
\lsmtestCatName{XY}{script}

All 17 real-world arguments classified correctly.
\end{document}
